% chktex-file 8
% chktex-file 12
% chktex-file 18
\beginsong{Okoř}

\beginverse{}
\ch{D}{}{}{}Na Okoř je cesta, jako žádná ze sta,
\ch{A7}{}{}{}roubená je stroma\ch{D}{}{}{}ma.
\ch{D}{}{}{}Když jdu po ní v létě, samoten na světě,
\ch{A7}{}{}{}sotva pletu noham\ch{D}{}{}{}a.
\ch{G}{}{}{}Na konci té cesty \ch{D}{}{}{}trnité,
\ch{E7}{}{}{}stojí krčma jako hrad, jako \ch{A7}{}{}{}hrad.
\ch{D}{}{}{}Tam zapadli trampi, hladoví a sešlí,
\ch{A7}{}{}{}začli sobě noto\ch{D}{}{}{}vat.
\endverse{}

\beginchorus{}
\ch{D}{}{}{}Na hradě Okoři \ch{A}{}{}{}světla už nehoří,
\ch{D}{}{}{}bílá paní \ch{A}{}{}{}šla už dávno\ch{D}{}{}{} spát.
Ta měla ve zvyku, \ch{A}{}{}{}podle svého budíku,
\ch{D}{}{}{}o půlnoci \ch{A}{}{}{}chodit straš\ch{D}{}{}{}ívat.
\ch{C}{}{}{}Od těch dob, co jsou tam \ch{G}{}{}{}trampové,
\ch{A7}{}{}{}nesmí z hradu \ch{D}{}{}{}pryč, z hradu pryč.
\ch{D}{}{}{}O půlnoci v podhradí \ch{A}{}{}{}se šerifem dovádí,
\ch{D}{}{}{}on jí sebr\ch{A}{}{}{}al od komnat\ch{D}{}{}{}y klíč.
\endchorus{}

\beginverse{}
Jednoho dne z rána roznesla se zpráva,
že byl Okoř vykraden.
Nikdo neví dodnes, kdo to tenkrát odnes,
nikdo nebyl dopaden.
Šerif hrál celou noc mariáš
s bílou paní v kostinci, v kostinci.
Místo aby hlídal, zuřivě ji líbal,
dostal z toho zimnici.
\endverse{}

\beginchorus{}
\phantom{}
\endchorus{}

\endsong{}
